% ---------------------------------------------------------------------------
% Experiments
% ---------------------------------------------------------------------------
\section{Experiments}
\label{sec:experiments}

We organize the evaluation into five experiments. The first two measure how far AHL
can improve against the human field; the third compares it with reinforcement learning
on a standard game benchmark; the final two isolate the role of opponent selection and
history. Unless stated otherwise, all runs use the same referee, frozen evaluation
pool, and append-only event schema described in Section~\ref{sec:arena-eval}.

\subsection{Fixed-budget performance against human programs}
\label{sec:exp-human}

Under a common revision budget, we run several coding-agent models on every game in
\benchname. For each model and game, we report the initial (raw) Elo, the Elo after
the fixed number of AHL iterations, and the gain. We also place the final policy in
the archived human leaderboard and report its rank and percentile. This experiment
measures cross-model performance at a matched budget rather than the benefit of a
single particularly long run.

\subsection{Long-horizon heuristic evolution}
\label{sec:exp-long}

We select one representative model and run it for a substantially longer horizon,
recording Elo after every revision on a diverse subset of games. The resulting curves
show the initial policy quality, the gains from continued maintenance, the point at
which AHL reaches the human reference, and whether improvement eventually plateaus.
We additionally relate the plateau to the maintained program's description length.

\subsection{AHL versus RL on Go or Chess}
\label{sec:exp-rl}

We compare AHL with established RL baselines on either Go or Chess, using the same
referee-facing interface and matched interaction budgets. The primary plot uses the
number of observed trajectories as the horizontal axis and Elo as the vertical axis;
we report both the final matched-budget Elo and the number of trajectories required to
reach a fixed Elo level. This isolates the sample-efficiency regime in which the
decision-level feedback of an editable program may compensate for the scaling
advantage of a learned policy.

\subsection{Ablation: opponent selection and pool size}
\label{sec:exp-opponent-ablation}

Keeping the model, game, initial policy, and total budget fixed, we vary how AHL
selects the policies that participate in each round and how many policies it samples.
We compare random sampling, agent-selected opponents, easy-to-hard scheduling, and
rank-based selection, together with multiple opponent-set sizes. We report final Elo,
Elo gain, and regression rate to determine whether improvements come from the identity
of the opponents, the diversity of the pool, or simply more matches per round.

\subsection{Ablation: the role of history}
\label{sec:exp-history-ablation}

We compare the standard history-aware agent with a history-free variant. The former
can inspect prior policies, replays, revisions, and outcomes; the latter receives only
the current policy and the latest feedback. All other components and budgets remain
fixed. The difference in final Elo and revision efficiency measures whether sustained
cross-round context is necessary for maintaining a policy rather than merely producing
independent edits.
